
        \documentclass[fleqn]{article}      
        \usepackage{amsmath}
        \usepackage{fontspec}
        \usepackage{kotex} % 한국어 지원  

        \begin{document}      
         객관식 질문

1. 적분의 기본 목적은 무엇인가요?
   a) 함수의 최대값을 찾기 위해  
   b) 함수의 면적을 계산하기 위해  
   c) 함수의 기울기를 계산하기 위해  
   d) 함수의 변화를 측정하기 위해  

2. 다음 중 정적분의 정의에 해당하는 것은?
   a) 함수의 기울기를 구하는 것  
   b) 함수의 값을 더하는 것  
   c) 함수의 면적을 구하는 것  
   d) 함수의 극한을 구하는 것  

3. 불연속 함수의 적분은 어떻게 처리하나요?
   a) 항상 0으로 계산한다  
   b) 적분을 계산할 수 없다  
   c) 불연속점을 제외하고 계산한다  
   d) 모든 점에서 계산한다  

4. 정적분의 기호는 무엇인가요?
   a) ∑  
   b) ∫  
   c) Δ  
   d) ∏  

5. 다음 중 적분의 결과가 음수가 될 수 있는 경우는 언제인가요?
   a) 함수가 항상 양수일 때  
   b) 함수가 음수인 구간이 있을 때  
   c) 함수가 0일 때  
   d) 모든 경우에서 항상 그렇다  

 주관식 질문

1. 적분이란 무엇인지 간단히 설명해 보세요.
2. 정적분을 사용하여 면적을 구하는 방법은 무엇인가요?
3. 적분과 미분의 관계에 대해 설명해 보세요.
4. 정적분의 기본정리에 대해 설명해 보세요.
5. 적분을 통해 구할 수 있는 실제 예시를 하나 들어보세요.
6. 함수 f(x) = x²의 정적분을 0에서 2까지 계산해 보세요.
7. 적분을 통해 평균값 정리를 설명해 보세요.
8. 함수의 불연속점에서의 적분 계산 방법을 설명해 보세요.
9. 정적분과 부정적분의 차이점은 무엇인가요?
10. 실생활에서 적분이 사용되는 예시를 설명해 보세요.

 쉬운 문제

1. 다음 함수의 정적분을 구하세요: f(x) = 2, x = 0에서 x = 3까지.
2. f(x) = x의 그래프 아래 면적을 x = 1에서 x = 4까지 구하세요.
3. f(x) = 3의 적분을 x = 2에서 x = 5까지 구하세요.
4. f(x) = x²의 정적분을 x = 1에서 x = 3까지 구해 보세요.
5. f(x) = 4의 정적분을 x = 0에서 x = 2까지 구하세요.

 중간 문제

1. f(x) = 2x + 1의 정적분을 x = 0에서 x = 2까지 구하세요.
2. f(x) = x³의 정적분을 x = 1에서 x = 2까지 구해 보세요.
3. f(x) = sin(x)의 정적분을 x = 0에서 x = π까지 구하세요.
4. f(x) = e^x의 정적분을 x = 0에서 x = 1까지 구하세요.
5. f(x) = 1/x의 정적분을 x = 1에서 x = 2까지 구해 보세요.

 어려운 문제

1. f(x) = x² + 2x + 1의 정적분을 x = -1에서 x = 1까지 구하세요.
2. f(x) = cos(x)의 정적분을 x = 0에서 x = π/2까지 구하세요.
3. f(x) = ln(x)의 정적분을 x = 1에서 x = e까지 구해 보세요.
4. f(x) = x^4의 정적분을 x = 0에서 x = 2까지 구해 보세요.
5. f(x) = 1/(x² + 1)의 정적분을 x = 0에서 x = 1까지 구하세요.

      
        \end{document} 
    